\documentclass[border=3pt]{standalone}
\usepackage{ctex}
\usepackage{tikz}
\usetikzlibrary{arrows.meta,positioning,shapes.geometric,calc}
\begin{document}
\begin{tikzpicture}[
  font=\small,
  box/.style={draw=#1!70!black, fill=#1!12, rounded corners=2pt, align=center,
              inner sep=4pt, minimum height=8mm, line width=0.8pt},
  arrow/.style={-{Stealth[length=2.5mm]}, line width=0.8pt, color=black!60}
]
\node[box=orange, text width=18mm] (rc) at (0,0) {背景字段 $x_c$\\[1pt]{\scriptsize (可缺失)}};
\node[box=green!60!black, text width=18mm] (ri) at (0,-1.7) {概要字段 $x_i$};
\node[box=blue, text width=20mm] (ec) at (2.8,0) {背景编码器\\[1pt]{\scriptsize $q_c(z\mid x_c)$}};
\node[box=blue, text width=20mm] (ei) at (2.8,-1.7) {概要编码器\\[1pt]{\scriptsize $q_i(z\mid x_i)$}};
\node[box=yellow!80!black, text width=10mm] (tc) at (5.4,0) {$\theta_c$};
\node[box=yellow!80!black, text width=10mm] (ti) at (5.4,-1.7) {$\theta_i$};
\node[box=cyan!70!black, text width=18mm] (fuse) at (8.1,-0.85) {门控融合\\[1pt]{\scriptsize $\theta^{\mathrm{fuse}}$}};
\node[box=violet, text width=20mm] (dec) at (10.9,-0.85) {共享解码器\\[1pt]{\scriptsize $\beta$}};
\node[box=purple!70!black, text width=24mm] (rec) at (13.7,-0.85) {跨字段词重构\\[1pt]{\scriptsize 背景词 $+$ 概要词}};
\node[box=red!70!black, text width=26mm] (js) at (8.1,-3.1) {JS$(\theta_c,\theta_i)$\\[1pt]{\scriptsize 一致性正则}};
\node[draw=black!70, fill=black!8, rounded corners=2pt, align=center, inner sep=4pt,
      minimum height=8mm, line width=0.8pt, text width=20mm] (loss) at (11.6,-3.1)
      {总损失 $\mathcal{L}$};
\draw[arrow] (rc) -- (ec);
\draw[arrow] (ri) -- (ei);
\draw[arrow] (ec) -- (tc);
\draw[arrow] (ei) -- (ti);
\draw[arrow] (tc.east) -- (fuse.west);
\draw[arrow] (ti.east) -- (fuse.west);
\draw[arrow] (fuse.east) -- (dec.west);
\draw[arrow] (dec.east) -- (rec.west);
\draw[arrow] (tc.east) -- ++(0.5,0) |- (js.north);
\draw[arrow] (ti.east) -- ++(1.1,0) |- (js.north);
\draw[arrow] (rec.south) |- (loss.east);
\draw[arrow, dashed] (js.east) -- (loss.west);
\node[font=\scriptsize, above=1mm of dec, text width=40mm, align=center]
      {两条流各自解码同一 $B$ 并重建全部存在字段（式 \eqref{eq:lrec}）};
\node[red!70!black, font=\scriptsize, above=1mm of rc, text width=26mm, align=center]
       {背景缺失时该流回退到先验};
\end{tikzpicture}
\end{document}
